\section{De Chat a Agent: un cambio de paradigma}

Al inicio de la entrevista, 2026 se define como el segundo acto de la guerra de los grandes modelos. El primer acto es Chat: el usuario y el modelo interactúan mediante preguntas y respuestas, y la competencia gira en torno a la escala del preentrenamiento, la alineación con instrucciones y la experiencia con contextos cortos. El segundo acto es Agent: el modelo entra en un entorno, invoca herramientas y ejecuta tareas de largo alcance, y la competencia se desplaza hacia el postentrenamiento, la RL infra, la retroalimentación del entorno, el costo y la confiabilidad.

\begin{figure}[H]
\centering
\includegraphics[width=0.94\textwidth]{chat-to-agent-paradigm.png}
\caption{Cambio de paradigma: de un predominio de Chat/Pre-train a un predominio de Agent/Post-train.}
\end{figure}

\begin{knowledgebox}{Lectura de la figura: cómo interpretar esta imagen}
A la izquierda, la etapa Chat pone el acento en las preguntas y respuestas, la recuperación de información y la generación, y el peso de los recursos recae en los datos de preentrenamiento y el modelo base; a la derecha, la etapa Agent pone el acento en las herramientas, el entorno y las tareas de largo alcance, y el peso de los recursos se traslada al rollout, la RL infra, la evaluación y los marcos. La flecha de la figura no dice que el preentrenamiento haya dejado de importar, sino que el foco de la competencia se ha desplazado.
\end{knowledgebox}

\subsection{El papel detonante de OpenClaw}

Al principio, Luo Fuli (罗福莉) veía OpenClaw como Claude Code con una interfaz de usuario añadida y rechazaba por instinto su envoltorio de mercadotecnia. Pero, tras usarlo en serio, lo definió como un marco de Agent que marca época, porque permite que el marco compense las carencias del modelo y que los miembros comunes de un equipo participen, a través del marco, en «elevar el nivel de inteligencia». La importancia de OpenClaw no está en una interfaz concreta, sino en que convierte el marco de Agent en un paradigma de investigación.

\begin{importantbox}{La esencia de OpenClaw}
OpenClaw organiza el modelo, las herramientas, el entorno de la tarea y la imaginación humana dentro de un marco iterable. No es solo una «innovación en la forma del producto», sino que hace explícita la pregunta de investigación de la era del postentrenamiento: cómo lograr que el marco de Agent y el modelo se potencien mutuamente.
\end{importantbox}

\subsection{Por qué 2026 es el año de la transformación de la productividad}

Luo Fuli considera que los desarrolladores chinos reaccionan con más fuerza ante herramientas del tipo de OpenClaw, por un lado porque su necesidad de mejorar la eficiencia es más apremiante, y por otro porque en China hay una gran cantidad de modelos baratos y fáciles de usar. Si una tarea compleja de Agent cuesta 10 yuanes de API pero sustituye un trabajo humano valorado en 1,000 yuanes, el incentivo para adoptarla es muy fuerte. Las condiciones previas de la revolución de productividad de los Agent son un marco lo bastante bueno, modelos lo bastante baratos y un retorno de la tarea lo bastante claro.

\begin{warningbox}{No confundir popularidad con madurez}
Que un marco se vuelva popular puede indicar que respondió a una necesidad de productividad, pero no significa que ya tenga madurez de nivel industrial. La madurez real exige cumplir a la vez con la tasa de finalización de extremo a extremo, la eficiencia de costos, la velocidad, la confiabilidad y los límites de seguridad.
\end{warningbox}

\subsection{Resumen de la sección}

La etapa Chat muestra la inteligencia del modelo; la etapa Agent exige que el modelo entre en el entorno de la tarea. El impacto de OpenClaw consiste en que muchas personas percibieron por primera vez que el marco, las herramientas y las formas de uso colectivo pueden cambiar con rapidez el límite superior visible de la capacidad del modelo.
